% page 525 of the facsimile (image p-524 of the scan)
% block 165.1 x 237.7 mm ; left margin 36.9 mm
\pgc{15.240mm}{0.000mm}{
\l{\cel{52}517}
\l{}
\l{\sou{selakto}. Parto seroza qua su separas de lakto kande ica koagulas.}
\l{}
\l{\sou{selektar}. (\sou{trans.})\cel{2}I. Prenar prefere. - II. Rezolvar por un}
\l{\cel{3}ek du o plura kozi. - DEFRS.}
\l{}
\l{\sou{selenio}. (\sou{kemio)}. Korpo simpla, metaloido vicina ye sulfo.}
\l{\cel{3}Atom-nombro : 34. Atom-pezo : 79,2. -\cel{1}\sou{Simb. kemiala}\cel{1}: Se.}
\l{\cel{3}- DEFIRS.}
\l{}
\l{\sou{selenografio.}\cel{2}Deskripto astronomiala di Luno. - DEFIRS.}
\l{}
\l{\sou{semar}. (\sou{trans.}) Disjetar, disenterigar, distope, en kultivo-sulo}
\l{\cel{3}preparata (la semini di ta o ca planto), o disjetante, o per}
\l{\cel{3}semo-mashino. - (\sou{anke metaf}.)- FIRS.}
\l{}
\l{\sou{semaforo}.\cel{2}Aparato telegrafala, instalita vicine ye portuo,}
\l{\cel{3}por signalar laadveno e la arivo di la navi, lia manovri}
\l{\cel{3}vide de la litoro di portuo, e por korespondar kun li.-}
\l{\cel{3}DEFIRS.}
\l{}
\l{\sou{semano}. Periodo de sep dii, de sundio til saturdio (inklu-}
\l{\cel{3}zite). - FIS.}
\l{}
\l{\sou{semantiko.}\cel{2}Studio di la senco di la vorti e di lia varii -}
\l{\cel{3}cienco pri la \textquotedbl{}vivo di la vorti\textquotedbl{}. - DEFIRS.}
\l{}
\l{\sou{semblar}. (\sou{netrans.}). Aspektar quale se lu esus (ma ne esas}
\l{\cel{3}reale). -EFI.}
\l{}
\l{\sou{semestro}. Periodo de sis monati intersequanta. - DEFIRS.}
\l{}
\l{\sou{semino}. (\sou{bot.)}\cel{2}Parto di frukto (grano, kerno, kerneto) quan onu}
\l{\cel{3}semas por produktar la planto quan lu kontenas jerme. - EFIS.}
\l{}
\l{\sou{seminario}. (\sou{relig}io). Domo religiala ube onu preparas, en sin-}
\l{\cel{3}gla diocezo (katolika) la kleriki aprentisa, qui aspiras a}
\l{\cel{3}nominesor kleriki. - DEFIRS.}
\l{}
\l{\sou{semlo.}\cel{1}Pano-bloko mikra, generale bulatra, ek frumento delikata.}
\l{}
\l{\sou{semolo}. Gruelo de frumento bakita, e granetigita per pisto o}
\l{\cel{3}trituro, quan onu uzas por facar supi ed entremesi. - EFIS.}
\l{}
\l{\sou{sempervivo.}\cel{1}(\sou{bot.}) Planto dika-folia, herbatra, ek la familio}
\l{\cel{3}\textquotedbl{}kresulacei\textquotedbl{}. - L.\cel{1}\sou{sempervivum tectorum.}\cel{2}L.}
\l{}
\l{\sou{sempre.}\cel{1}I.\textquotedbl{}En la tempo futura, sen mem nur un ecepto-kazo\textquotedbl{} -}
\l{\cel{3}- II. \textquotedbl{}En la tempo pasinta, sen mem nur un ecepto-kazo\textquotedbl{}.}
\l{\cel{3}- III.\textquotedbl{}Ye irga epoki\textquotedbl{}. - EFIS.}
\l{}
\l{sen. Preppziciono qua indikas la absenteso, la manko, la indijo}
\l{\cel{3}di persono, di kozo, di la eso-maniero o di la ago-maniero}
\l{\cel{3}quan lu guvernas. - FIS.}
}
